% Incorporación conservativa: selección aditiva e instrumento térmico.
\section{Selección aditiva e instrumento de la lectura térmica}
\label{sec:ii-aditiva-instrumento}

La referencia espectral y la acción construidas en las secciones anteriores
permiten precisar qué se cuenta al medir información y energía. El coeficiente
térmico aparece como una intensidad relativa entre dos sucesos del mismo
registro. Su función informacional queda caracterizada por la aditividad, y
ambos sucesos admiten un aparato explícito que conserva el estado completo.
Se desarrolla así la regla marcada de
\ref{thm:ii-ter27-transducer}, con su realización, su selección dentro de una
clase de lectores y sus condiciones de transporte.

La composición mantiene el antecedente APP--TRIT--TPK: hojas aditiva y
multiplicativa, régimen, orientación, residuo, cociente, frontera, ruta y
memoria. La prolongación
$w_6\to w_{12}\to w_{18}\to w_{24}\to w_{30}\to R_{36}\to G_9$
y el orden $U_t\to\Gamma_9\to K_{\rm ph}$ preceden a esta lectura
espectral. El retorno de fase prolonga el registro; las marcas observadas
no sustituyen ese registro ni las dos proyecciones de la estructura discreta
conjunta del continuo. Se reciben la razón $r=1/10$, los rangos
$(1,380,649)$ y su asignación a los grados $23+3n$ con las condiciones de
\ref{subsec:ii-ter27-reference}. Esta asignación se conserva como parte
explícita de la realización, junto con el Hamiltoniano y su origen energético.

% BEGIN THERMAL_R03_01
\subsection{Sucesos espectrales y medida relativa}
\label{subsec:ii-tins-eventos}

Sea $\mathcal K=\ell^2(\mathbb N_0)\otimes\mathbb C^M$, con $M=649$,
y consideremos la densidad normalizada
\begin{equation}
 \Sigma=\rho_r\otimes I_M/M,\qquad
 \rho_r=(1-r)\sum_{j\ge0}r^j|j\rangle\langle j|.
 \label{eq:ii-tins-sigma}
\end{equation}
Para $n=0,1,2$, sean $j_n=23+3n$ y $Q_n$ proyectores de memoria de rango
$d_n$, donde $(d_0,d_1,d_2)=(1,380,649)$. Fijemos un vector unitario $\chi_B$.
Las letras $A,B,C$ designan aquí sucesos de lectura, no el área geométrica:
\begin{align}
 P_A&=\sum_{n=0}^2|j_n\rangle\langle j_n|\otimes Q_n,&
 P_B&=|0\rangle\langle0|\otimes|\chi_B\rangle\langle \chi_B|,
 \nonumber\\
 P_C&=I-P_A-P_B,& \mu(P)&=\Tr(\Sigma P).
 \label{eq:ii-tins-marks}
\end{align}
Los tres proyectores son ortogonales, suman la identidad y conmutan con
$\Sigma$. El suceso $B$ fija una referencia de rango uno; $C$ retiene todo
el complemento de los sectores marcados.

\Needspace{12\baselineskip}
\begin{samepage}
\begin{proposition}[Intensidad relativa y compresión]
\label{prop:ii-tins-relative}
Escribiendo $p_A=\mu(P_A)$ y $p_B=\mu(P_B)$, se obtiene
\begin{equation}
 p_B=\frac{1-r}{M},\qquad
 p_A=\frac{1-r}{M}\sum_{n=0}^2d_nr^{23+3n},\qquad
 \frac{p_A}{p_B}=b_*=
 \frac{1380649}{10^{29}}=1.380649\times10^{-23}.
 \label{eq:ii-tins-relative}
\end{equation}
Si $\iota:\mathcal D\to\mathcal K$ incluye cada sector del portador
en el correspondiente grado $j_n$ y en $\operatorname{ran}Q_n$, entonces
\begin{equation}
 p_B^{-1}\iota^*\Sigma\iota
 =\bigoplus_{n=0}^2r^{23+3n}I_{d_n}=\eta.
 \label{eq:ii-tins-compression}
\end{equation}
\end{proposition}
\end{samepage}
\begin{proof}
En el grado $j$, la densidad actúa como $(1-r)r^j/M$ por la identidad
de memoria. La traza sobre un subespacio de rango $d$ multiplica ese
escalar por $d$. Aplicarlo al grado cero y a los tres grados marcados
da las dos probabilidades; su cociente cancela el factor común.
Como $r=10^{-1}$,
$\sum_nd_nr^{23+3n}=10^{-23}(1+380\,10^{-3}+649\,10^{-6})$,
que da la fracción indicada. La misma acción escalar, restringida mediante
$\iota$, demuestra \eqref{eq:ii-tins-compression}.
\end{proof}

La intensidad $b_*$ es una medida relativa, no la probabilidad bruta $p_A$.
Usar como referencia todo el vacío $|0\rangle\langle0|\otimes I_M$
produciría $b_*/M$: cambia el suceso de referencia. Añadir un factor auxiliar
independiente y extender ambos proyectores por su identidad conserva, en
cambio, las dos probabilidades y el cociente.
% END THERMAL_R03_01

% BEGIN THERMAL_R03_02
\subsection{Aparato unitario, recuperación y dominio energético}
\label{subsec:ii-tins-aparato}

En un puntero de dimensión tres, con base $|A\rangle,|B\rangle,|C\rangle$,
definamos
\begin{equation}
 V\psi=\sum_{y=A,B,C}P_y\psi\otimes|y\rangle.
 \label{eq:ii-tins-isometry}
\end{equation}
Para realizarlo mediante un unitario, sean $R_C=I$, $R_A$ el intercambio
de $|C\rangle$ con $|A\rangle$ y $R_B$ el intercambio de $|C\rangle$
con $|B\rangle$. El acoplamiento es
\begin{equation}
 \mathbb U=\sum_{y=A,B,C}P_y\otimes R_y,
 \qquad \mathbb U(\psi\otimes|C\rangle)=V\psi.
 \label{eq:ii-tins-unitary}
\end{equation}

\begin{proposition}[Registro conservativo de las tres alternativas]
\label{prop:ii-tins-instrument}
Se cumple $\mathbb U^*\mathbb U=\mathbb U\mathbb U^*=I$ y $V^*V=I$.
Para cualquier densidad $\xi$, el estado conjunto $\widehat\xi=V\xi V^*$
permite la recuperación exacta $V^*\widehat\xi V=\xi$. La lectura del
puntero define el instrumento
\begin{equation}
 \mathcal I_y(\xi)=P_y\xi P_y,\qquad
 p_y=\Tr(\xi P_y),\qquad
 \xi_y=\mathcal I_y(\xi)/p_y\quad(p_y>0).
 \label{eq:ii-tins-instrument}
\end{equation}
Su canal no selectivo preserva la traza y deja $\Sigma$ invariante.
\end{proposition}
\begin{proof}
Los productos cruzados de los $P_y$ se anulan. Por tanto,
$\mathbb U^*\mathbb U=\sum_yP_y\otimes R_y^*R_y=I$, y análogamente
para el otro producto. La ortonormalidad del puntero da $V^*V=\sum_yP_y=I$;
aplicarla a ambos lados de $\widehat\xi$ demuestra la recuperación.
Comprimir el puntero a $|y\rangle$ da \eqref{eq:ii-tins-instrument}.
Además, $\sum_y\Tr(P_y\xi P_y)=\Tr\xi$. Como cada $P_y$ conmuta con
$\Sigma$, también $\sum_yP_y\Sigma P_y=\Sigma$.
\end{proof}

Para un estado con coherencias entre marcas, olvidar el puntero puede
cambiarlas. La recuperación anterior utiliza el estado conjunto, que las
conserva. Leer un resultado y conservar el aparato completo son operaciones
distintas.

Sea $J|j\rangle=j|j\rangle$ y sea $H_{\rm sp}=f(J)\otimes I_M$, con
$f:\mathbb N_0\to\mathbb R$. Su dominio autoadjunto es
\[
 \mathcal D(H_{\rm sp})=
 \left\{\psi:\sum_{j,a}|f(j)|^2|\psi_{j,a}|^2<\infty\right\}.
\]
Los proyectores de las marcas reducen este dominio y conmutan con las
proyecciones espectrales de $H_{\rm sp}$. En él se verifica
\begin{equation}
 (H_{\rm sp}\otimes I_3)V=VH_{\rm sp},\qquad
 V^*(H_{\rm sp}\otimes I_3)V=H_{\rm sp}.
 \label{eq:ii-tins-energy}
\end{equation}
En efecto, aplicar $H_{\rm sp}$ a cada sumando de
\eqref{eq:ii-tins-isometry} permite permutarlo con $P_y$; la norma
cuadrática de la suma es la suma de las normas cuadráticas de sus sectores.
La segunda igualdad sigue de $V^*V=I$. Se conservan así las expectativas
energéticas cuando existen, incluida la realización
$f(j)=\epsilon_a(j+1/2)$. Un Hamiltoniano propio del puntero requiere su
especificación: si sus tres niveles son degenerados, el acoplamiento conserva
también la energía total. Esta identidad no asigna un coste nulo a la
preparación o al reinicio físico del aparato.
% END THERMAL_R03_02

% BEGIN THERMAL_R03_03
\subsection{Transporte trítico con residuo y acarreo}
\label{subsec:ii-tins-tritico}

La división $j=3w+b$, $0\le b<3$, define el unitario
$U_3|j\rangle=|w,b\rangle$. Sea $W=U_3\otimes I_M$. Las marcas transportadas son
\begin{align}
 P_A'&=\sum_{n=0}^2|7+n,2\rangle\langle7+n,2|\otimes Q_n,\nonumber\\
 P_B'&=|0,0\rangle\langle0,0|\otimes|\chi_B\rangle\langle \chi_B|,
 \qquad P_C'=I-P_A'-P_B'.
 \label{eq:ii-tins-tritic-marks}
\end{align}

\begin{proposition}[Naturalidad del aparato]
\label{prop:ii-tins-tritic}
Si $V'$ es el aparato de las marcas transportadas, entonces
\begin{equation}
 V'W=(W\otimes I_3)V,\qquad
 \mathcal I_y'(W\xi W^*)=W\mathcal I_y(\xi)W^*.
 \label{eq:ii-tins-square}
\end{equation}
Las probabilidades y los estados relativos se transportan sin pérdida.
\end{proposition}
\begin{proof}
La división euclídea da $P_y'W=WP_y$ para cada $y$. Sustituir estas
tres identidades en la suma que define $V'$ prueba el primer cuadrado;
multiplicar por los adjuntos prueba el segundo. La invariancia de la traza
bajo el unitario demuestra la afirmación probabilística, y dividir por
la misma probabilidad positiva demuestra la de los estados relativos.
\end{proof}

La densidad transportada conserva separadamente cociente y residuo:
\begin{equation}
 W\Sigma W^*=\rho_{r^3}\otimes\sigma_r\otimes I_M/M,
 \qquad \sigma_r=\frac{\operatorname{diag}(1,r,r^2)}{1+r+r^2}.
 \label{eq:ii-tins-tritic-state}
\end{equation}
En efecto, $(1-r)r^{3w+b}=(1-r^3)(r^3)^w r^b/(1+r+r^2)$.
Para el desplazamiento unilateral $S|j\rangle=|j+1\rangle$ se obtiene
\begin{equation}
 U_3SU_3^*=I\otimes\bigl(|1\rangle\langle0|+|2\rangle\langle1|\bigr)
 +S_w\otimes|0\rangle\langle2|.
 \label{eq:ii-tins-carry}
\end{equation}
Los casos $b=0,1$ incrementan el residuo; $b=2$ lo devuelve a cero e
incrementa $w$. Esta prueba conserva el acarreo. El desplazamiento puede
cambiar la marca: \eqref{eq:ii-tins-square} expresa independencia del orden
entre observación y cambio de representación, no inmovilidad bajo la dinámica.
% END THERMAL_R03_03

% BEGIN THERMAL_R03_04
\subsection{Inclusiones equivalentes y ampliación de memoria}
\label{subsec:ii-tins-inclusiones}

Sean $\widetilde Q_n$ otros proyectores de los mismos rangos y
$\widetilde \chi_B$ otro vector unitario. En cada grado $j_n$ existe un unitario
$R_{j_n}$ que lleva $\operatorname{ran}Q_n$ a
$\operatorname{ran}\widetilde Q_n$: basta enviar bases ortonormales de esos
subespacios y de sus complementos. Elijamos igualmente $R_0\chi_B=\widetilde \chi_B$
y completemos con unitarios en los restantes grados. Entonces
$R=\bigoplus_{j\ge0}R_j$ es unitario, conserva $\Sigma$ y conmuta con
$H_{\rm sp}$. Las identidades
\begin{equation}
 \widetilde P_yR=RP_y,\qquad
 \widetilde V R=(R\otimes I_3)V
 \label{eq:ii-tins-inclusion}
\end{equation}
se verifican bloque a bloque. Por conjugación y ciclicidad de la traza,
conservan probabilidades, energía espectral y $b_*$. No se requiere que un
único unitario de memoria independiente del grado realice simultáneamente
todas las inclusiones. Cualquier operador adicional se transporta por el
mismo $R$; esta conjugación no afirma que su matriz permanezca inalterada.

Para una ampliación, sea $L:\mathbb C^M\to\mathbb C^{M'}$ una isometría
y sea $J_L=I\otimes L$. Definamos
$P_A^+=J_LP_AJ_L^*$, $P_B^+=J_LP_BJ_L^*$ y
$P_C^+=I-P_A^+-P_B^+$. Como $J_L^*J_L=I$, para las tres marcas
$P_y^+J_L=J_LP_y$. En consecuencia,
\begin{equation}
 V^+J_L=(J_L\otimes I_3)V,
 \qquad \Sigma_L=J_L\Sigma J_L^*=\rho_r\otimes LL^*/M.
 \label{eq:ii-tins-extension}
\end{equation}
El transporte isométrico conserva las probabilidades originales.
Si, además, se prepara la densidad uniforme de todo el ambiente ampliado,
$\Sigma^+=\rho_r\otimes I_{M'}/M'$, la preparación es distinta y
\begin{equation}
 p_A^+=\frac{M}{M'}p_A,\qquad
 p_B^+=\frac{M}{M'}p_B,\qquad p_A^+/p_B^+=b_*.
 \label{eq:ii-tins-uniform-extension}
\end{equation}
La prueba vuelve a ser la traza del escalar de cada grado sobre los mismos
rangos. En $A$ y $B$, normalizar cancela $M/M'$ y produce las imágenes
isométricas de los estados relativos originales; el complemento recibe el
peso añadido. Así se distinguen conservación del aparato, igualdad de la
intensidad e igualdad del estado completo.
% END THERMAL_R03_04

% BEGIN THERMAL_R03_05
\subsection{Preparación completa, entropía y condicionamiento}
\label{subsec:ii-tins-condicionamiento}

Sea $\rho$ un estado de un sistema de ensayo finito, $H=H^*$ su Hamiltoniano
heredado y $\rho_0$ una preparación pura fija. Con $P=P_A$, definamos
\begin{equation}
 \Omega_\rho=(P\Sigma P)\otimes\rho
 +((I-P)\Sigma(I-P))\otimes\rho_0,
 \qquad \Omega_0=\Sigma\otimes\rho_0.
 \label{eq:ii-tins-preparation}
\end{equation}
Esta preparación conserva todo el complemento espectral. Es distinta del
producto auxiliar de \eqref{eq:ii-ter27-flag}, aunque realiza los mismos dos
funcionales sobre el ensayo.

\begin{theorem}[Información por referencia y energía condicionada]
\label{thm:ii-tins-readings}
La traza de $\Omega_\rho$ vale uno, y la traza parcial sobre el espacio
de ensayo devuelve $\Sigma$. Además,
\begin{align}
 \mathcal S_B(\rho)&:=
 \frac{s(\Omega_\rho)-s(\Omega_0)}{p_B}=b_*s(\rho),
 \label{eq:ii-tins-information}\\
 \rho_A&:=\frac{\Tr_{\mathcal K}[(P\otimes I)\Omega_\rho(P\otimes I)]}{p_A}
 =\rho,\qquad
 \mathcal E_A(\rho)=\Tr(\rho_AH)=\Tr(\rho H).
 \label{eq:ii-tins-conditioned}
\end{align}
\end{theorem}
\begin{proof}
Los sumandos positivos tienen trazas $p_A$ y $1-p_A$. Como $[P,\Sigma]=0$,
la traza parcial sobre el ensayo es
$P\Sigma P+(I-P)\Sigma(I-P)=\Sigma$. Sean $\lambda_i$ los autovalores
del bloque marcado y $p_j$ los de $\rho$. Sus productos sustituyen a
$\lambda_i$; el incremento entrópico de ese bloque es
\[
 -\sum_{i,j}\lambda_ip_j\ln(\lambda_ip_j)
 +\sum_i\lambda_i\ln\lambda_i=p_As(\rho).
\]
El complemento no cambia. Las sumas son legítimas, pues
$s(\Sigma)=-\ln(1-r)-r\ln r/(1-r)+\ln M<\infty$ y el ensayo es finito.
Dividir por $p_B$ demuestra \eqref{eq:ii-tins-information}. Comprimir a
$P\otimes I$ y tomar la traza parcial da $p_A\rho$; dividir por $p_A>0$
demuestra \eqref{eq:ii-tins-conditioned} y su expectativa energética.
\end{proof}

El condicionamiento conserva la operación observacional del corpus. Si
$U_a$ acopla el estado con el registro inicial $m_0$, la salida y su fibra son
\[
 \operatorname{Out}_a(x)=q_a(\operatorname{pr}_M U_a(x,m_0)),\qquad
 \Omega_{a,y}=\{x:\operatorname{Out}_a(x)=y\}.
\]
Para $\mu_a(\Omega_{a,y})>0$ se restringe la medida a esa fibra:
\begin{equation}
 \mu_{a,y}(D)=\frac{\mu_a(D\cap\Omega_{a,y})}{\mu_a(\Omega_{a,y})}.
 \label{eq:ii-tins-fiber}
\end{equation}
Cuando esos sucesos pertenecen a una misma medida espectral proyectiva y
$\mu_a(D)=\Tr(\omega P_D)$, se tiene
$P_DP_y=P_{D\cap\Omega_{a,y}}$. La ciclicidad de la traza prueba
\[
 \Tr\!\left(\frac{P_y\omega P_y}{\Tr(\omega P_y)}P_D\right)
 =\mu_{a,y}(D).
\]
La igualdad se extiende a observables acotados de esa álgebra por sumas
espectrales y aproximación uniforme. Aplicada al suceso $A$ de
\eqref{eq:ii-tins-preparation}, da exactamente la energía media
condicionada anterior. No requiere una hipótesis dinámica de ergodicidad.
Una purificación adicional, cuando exista, se conserva en su propio factor;
la entropía aquí calculada pertenece al álgebra observable especificada.
% END THERMAL_R03_05

% BEGIN THERMAL_R03_06
\subsection{Caracterización por aditividad y refinamiento}
\label{subsec:ii-tins-aditividad}

La evaluación anterior admite una caracterización que fija su peso sin
tomar $b_*$ como coeficiente postulado. Fijado el estado de ensayo $\rho$,
consideremos una contribución $F_\rho(x)$ para una alternativa de peso
recibido $x\in[0,1]$. En la lectura relativa, ese peso se mide respecto de
la referencia $B$; sólo se consideran alternativas del dominio indicado.
Exigimos: dependencia del peso y no del nombre de la alternativa; aditividad
$F_\rho(x+y)=F_\rho(x)+F_\rho(y)$ para $x+y\le1$; positividad; y
normalización $F_\rho(1)=s(\rho)$. El dominio debe realizar esos repartos,
o los refinamientos ordenados que se describen a continuación.

\begin{theorem}[Unicidad del peso informacional aditivo]
\label{thm:ii-tins-additive}
En el dominio escalar completo, las cuatro condiciones anteriores implican
$F_\rho(x)=xs(\rho)$. La misma conclusión vale para una realización de
alternativas con refinamientos equiponderados de pesos $3^{-n}$ y, para
cada marca $P$, aproximaciones $P_n^-\le P\le P_n^+$ cuyos pesos convergen
al de $P$ y cuyas contribuciones pertenecen a esos refinamientos.
\end{theorem}
\begin{proof}
La aditividad da $F_\rho(0)=0$. Si $x\le y$, entonces
$F_\rho(y)-F_\rho(x)=F_\rho(y-x)\ge0$, de modo que el lector es monótono.
La partición de la unidad en $3^n$ alternativas iguales da
$F_\rho(3^{-n})=3^{-n}s(\rho)$, y reunir $k$ de ellas da
$F_\rho(k3^{-n})=k3^{-n}s(\rho)$. Para $0\le x<1$ tomamos
\[
 \ell_n=3^{-n}\lfloor3^nx\rfloor,\qquad u_n=\ell_n+3^{-n}.
\]
Son cotas de $x$ dentro de $[0,1]$. La monotonía implica
$\ell_ns(\rho)\le F_\rho(x)\le u_ns(\rho)$; su diferencia
$3^{-n}s(\rho)$ tiende a cero. Esto prueba la fórmula, incluido
$s(\rho)=0$; $x=1$ es la normalización. Para marcas ordenadas se repite
el argumento con $P_n^-$ y $P_n^+$, pues la aditividad positiva implica
isotonicidad. No se presupone continuidad diferenciable del lector.
\end{proof}

En particular, la contribución de $A$ es $b_*s(\rho)$ cuando se conserva
su medida relativa. La caracterización se aplica también a la marca
normalizada de \eqref{eq:ii-ter27-flag}. Allí el funcional
$-\Tr[(P\otimes I)(\widehat\eta\otimes\rho)(I\otimes\ln\rho)]$
satisface las condiciones por linealidad de la traza y asigna $s(\rho)$
a la marca cierta.

El refinamiento requiere una realización efectiva: el proyector $P_B$ de
rango uno no contiene tres subproyectores no nulos equiponderados. Una
ampliación auxiliar puede realizar refinamientos, pero debe transportar
la medida y la lectura. Asimismo, tres descendientes de un cilindro no
implican por su cardinalidad pesos iguales. Para un peso racional basta
una partición equiponderada del denominador correspondiente, si existe en
la realización; la prueba ternaria no presupone ese denominador particular.

Las hipótesis separan exactamente las libertades descartadas. El lector
$f(x)=x^2$ respeta productos pero da
$f(x+y)-f(x)-f(y)=2xy$; con $x=1/3$, $y=1/5$, el defecto es $2/15$.
Sin normalización, $f(x)=cx$ conserva positividad y aditividad. Sin
identificar la medida recibida, $(1/3,2/3)$ y $(1/2,1/2)$ son dos medidas
positivas, aditivas y normalizadas diferentes sobre dos alternativas.
La entropía fija del registro de etiquetas se mantiene aparte: el teorema
no elimina el término de mezcla de una entropía conjunta.
% END THERMAL_R03_06

% BEGIN THERMAL_R03_07
\subsection{Realización por visitas entre retornos}
\label{subsec:ii-tins-retornos}

La intensidad relativa posee también una interpretación dinámica precisa.
Sea $(X,\mu,T)$ un sistema probabilístico invertible, ergódico y preservador
de medida, con sucesos $A,B$ de medidas $p_A,p_B>0$. Para $x\in B$, sean
\[
 R_B(x)=\min\{k\ge1:T^kx\in B\},\qquad
 C_A(x)=\sum_{k=0}^{R_B(x)-1}\mathbf1_A(T^kx).
\]
La excursión conserva su secuencia completa; $C_A$ es sólo uno de sus observables.

\begin{theorem}[Recompensa por retorno a una referencia]
\label{thm:ii-tins-return}
Con $\mu_B=\mu|_B/p_B$, para todo $f\in L^1(\mu)$,
\begin{equation}
 \mathbb E_{\mu_B}\!\left[\sum_{k=0}^{R_B-1}f\circ T^k\right]
 =\frac1{p_B}\int_Xf\,d\mu.
 \label{eq:ii-tins-return}
\end{equation}
En particular, $\mathbb E_{\mu_B}C_A=p_A/p_B=b_*$.
\end{theorem}
\begin{proof}
Sean $B_k=\{x\in B:R_B(x)>k\}$. Recurrencia, invertibilidad y
ergodicidad aseguran que los niveles $T^kB_k$, $k\ge0$, particionan $X$
salvo un conjunto nulo: cada punto tiene una última visita a $B$ en su
pasado. Para $f\ge0$, preservación de medida y convergencia monótona dan
\[
 \int_B\sum_{k=0}^{R_B-1}f(T^kx)\,d\mu
 =\sum_{k\ge0}\int_{T^kB_k}f\,d\mu=\int_Xf\,d\mu.
\]
Para $f\in L^1$, aplicar primero la identidad a $|f|$ justifica la
integrabilidad, y después se restan las partes positiva y negativa.
Dividir por $p_B$ prueba la fórmula; $f=\mathbf1_A$ da la consecuencia.
\end{proof}

En una cadena finita irreducible de matriz $P$ y distribución estacionaria
$\pi>0$, la prueba admite una forma algebraica. Partiendo del estado de
referencia $B$, sea $v_i$ el número esperado de visitas a $i$ antes del
primer retorno positivo. El retorno tiene esperanza finita; sumar las
transiciones hasta él muestra que llegadas y salidas tienen las mismas
esperanzas, luego $vP=v$. Como $v_B=1$, la unicidad de la distribución
estacionaria da $v_i=\pi_i/\pi_B$. Sumar sobre $A$ reproduce
\eqref{eq:ii-tins-return}. Agrupar la distribución espectral en $B$, los
tres sectores de $A$ y el complemento proporciona una realización finita
con el cociente exacto \eqref{eq:ii-tins-relative}.

Si cada visita a $A$ lleva un ensayo $\rho$, independiente de los restantes
condicionado al registro de la excursión, los ensayos forman
$\rho^{\otimes C_A}$. La aditividad tensorial da
\begin{align}
 \mathbb E_{\mu_B}s(\rho^{\otimes C_A})&=b_*s(\rho),\nonumber\\
 \frac{\mathbb E_{\mu_B}[C_A\Tr(\rho H)]}
      {\mathbb E_{\mu_B}C_A}&=\Tr(\rho H).
 \label{eq:ii-tins-visit-readings}
\end{align}
La segunda expresión es el cociente de energía total y visitas totales,
no una media de cocientes por excursión, que serían indefinidos cuando
$C_A=0$. La energía acumulada por retorno vale $b_*\Tr(\rho H)$.

La información por retorno y la energía por visita realizan, por tanto,
los dos denominadores de \eqref{eq:ii-tins-information}--
\eqref{eq:ii-tins-conditioned}. Si los ensayos están correlacionados,
se usa su entropía condicional conjunta. Un proceso no invertible admite
la formulación en su extensión bilateral estacionaria; sin ergodicidad,
el dominio debe restringirse al soporte barrido por la sección de retorno.
Estas hipótesis dinámicas no se necesitan para el teorema operatorio
\ref{thm:ii-tins-readings}. Ni una cadena estacionaria de prueba se identifica
por sí sola con la actualización TPK, ni $R_B$ con el período material
$108t_0$.
% END THERMAL_R03_07

% BEGIN THERMAL_R03_08
\subsection{Referencia común y correlaciones entre sistemas}
\label{subsec:ii-tins-comun}

Una sola referencia y $m$ sistemas independientes producen el registro
$\widehat\eta\otimes\rho_1\otimes\cdots\otimes\rho_m$, con una
única marca $P\otimes I$. Su peso sigue siendo $b_*$. Para
$H_{\rm tot}=\sum_jH_j$, con las extensiones tensoriales correspondientes,
\begin{equation}
 \mathcal S_{P,\rm tot}=b_*\sum_js(\rho_j),\qquad
 \mathcal E_{P,\rm tot}=\sum_j\Tr(\rho_jH_j).
 \label{eq:ii-tins-common}
\end{equation}
La prueba es la factorización de la traza y
$\ln(\rho_1\otimes\rho_2)=\ln\rho_1\otimes I+I\otimes\ln\rho_2$
en el soporte, iterada sobre los factores. Para un estado correlacionado
$\rho_{1\ldots m}$, la misma marca da $b_*s(\rho_{1\ldots m})$. En dos
sistemas,
\begin{equation}
 \mathcal S_{P,1}+\mathcal S_{P,2}-\mathcal S_{P,12}
 =b_*I(1:2)_\rho,
 \label{eq:ii-tins-mutual}
\end{equation}
por la definición de información mutua. Un Hamiltoniano con interacción
conserva también una sola referencia, aunque su estado de Gibbs no factorice.

Dos referencias independientes con marca conjunta $P_1\otimes P_2$
tienen peso $b_1b_2$; para referencias correlacionadas el peso es
$\Tr[\omega_{R_1R_2}(P_1\otimes P_2)]$ y no tiene por qué factorizar.
Copiar una etiqueta no crea independencia. En efecto, la isometría
\[
 V_Pv=Pv\otimes|1\rangle+(I-P)v\otimes|0\rangle
\]
satisface $V_P^*(P\otimes|1\rangle\langle1|)V_P=P$; la traza sobre
el estado transportado conserva $b_*$, no $b_*^2$. Más generalmente,
para una isometría $L$,
$\Tr[(L\omega L^*)(LPL^*)]=\Tr(\omega P)$, por $L^*L=I$.
Retornar conservando el mismo registro no eleva automáticamente su peso
a una potencia. Una nueva selección de supervivencia utiliza su propia marca.
% END THERMAL_R03_08

% BEGIN THERMAL_R03_09
\subsection{Selección termométrica y composición con acción y horizonte}
\label{subsec:ii-tins-seleccion}

Fijemos la referencia y un Hamiltoniano finito no escalar $H$. Para
$\rho_\beta=e^{-\beta H}/Z(\beta)$, $\beta>0$, el teorema aditivo y
la energía condicionada producen la pareja
$\mathcal S_B=b_*s(\rho_\beta)$, $\mathcal E_A=\Tr(\rho_\beta H)$.
Si $E=\Tr(\rho_\beta H)$, la suma espectral finita da
$E'=-\operatorname{Var}_{\rho_\beta}(H)<0$ y
$(\ln Z)'=-E$. Derivar $s=\beta E+\ln Z$ da $s'=\beta E'$ y, por tanto,
\begin{equation}
 \frac{d\mathcal S_B}{d\mathcal E_A}=b_*\beta,\qquad
 \Theta_{B,A}=\frac1{b_*\beta},\qquad B_*(u_\Theta)=b_*u_E.
 \label{eq:ii-tins-temperature}
\end{equation}
Las bases positivas satisfacen $u_S=u_E/u_\Theta$. La última ecuación
determina la única aplicación lineal positiva que envía
$\Theta_{B,A}u_\Theta$ a $\beta^{-1}u_E$. El balance
\[
 \Tr(\rho H)-b_*\Theta s(\rho)-
 \bigl[\Tr(\rho_\Theta^*H)-b_*\Theta s(\rho_\Theta^*)\bigr]
 =b_*\Theta D(\rho\Vert\rho_\Theta^*)
\]
con $\rho_\Theta^*\propto e^{-H/(b_*\Theta)}$ conserva el mínimo
único demostrado en \ref{cor:ii-ter27-variational}. La selección aditiva
caracteriza ahora el peso que entra en ese balance.

Con $\epsilon_a=h_a/(108t_0)=\hbar_ac/L_\alpha$,
$q_A=\gamma a_0L_\alpha^2$ y
$\mathcal T_{I/A}=\epsilon_a/q_A$, las secciones térmicas son
\begin{align}
 b_*\Theta_{{\rm clk},P,a}
 &=\frac{h_a}{108t_0\ln3}
 =\frac{\mathcal T_{I/A}q_A}{\ln3},\nonumber\\
 b_*\Theta_{H,P}(R)&=\frac{\hbar_ac}{2\pi R},&
 \frac{\Theta_{H,P}(R)}{\Theta_{{\rm clk},P,a}}
 &=\frac{L_\alpha}{R}\frac{\ln3}{2\pi}.
 \label{eq:ii-tins-action-area}
\end{align}
Son ecuaciones de coordenadas de $B_*$ en las bases transportadas.
La primera usa la sección $\beta_{{\rm clk},a}=\ln3/\epsilon_a$;
la segunda, la escala $\tau_R=\hbar_ac/(2\pi R)$. Dividirlas prueba
el cociente. La relación radial $G_a\hbar_a=c^3L_\alpha^2$ implica
\begin{equation}
 G_ab_*\Theta_{{\rm clk},P,a}\ln3=c^4L_\alpha.
 \label{eq:ii-tins-gravity}
\end{equation}
En el horizonte, $\mathcal S_B=b_*s$ y
$\Theta_{H,P}=\tau_R/b_*$ convierten el balance previo en
\[
 \delta E_R-\Theta_{H,P}\delta\mathcal S_B
 =\tau_RD(\sigma\Vert\rho_A).
\]
La composición mantiene la resolución de Barbero, la incidencia sectorial,
la acción y el origen energético: la cancelación de escalas conserva la
misma energía, no añade un parámetro metrológico al constructor.

La selección tiene un dominio definido. Si se condicionan ambos canales,
se obtiene $(E,s)$; si se cuentan ambos por referencia, $(b_*E,b_*s)$.
En los dos casos la derivada es $\beta$. La pareja aquí caracterizada es
$(E,b_*s)$ y utiliza energía por visita e información por referencia.
La conservación del archivo permite las tres observaciones, pero no las
identifica. Tampoco los productos energéticos seleccionan solos el peso:
para $p>0$, reemplazar $b_*$ por $b_*^p$ y dividir las temperaturas por
el mismo factor relativo conserva esos productos y la composición serial;
el lector cuadrático incumple, sin embargo, la aditividad demostrada arriba.

Quedan determinados el instrumento de las marcas, su recuperación y
transporte, la contribución informacional aditiva, la energía condicionada
y el transductor de esta realización. Identificarla con un instrumento
termométrico del constructor completo exige transportar hacia $A$ y $B$
sus sucesos, su medida y su par de normalizaciones, conservando el
refinamiento y la memoria. La existencia del aparato no sustituye esa
identificación particular. Si la referencia cambia con $\beta$, aparecen
además sus derivadas; si la proyección no conmuta con $\Sigma$, la preparación
\eqref{eq:ii-tins-preparation} conserva su diagonal por bloques y no
necesariamente la densidad original. Estas condiciones precisan el alcance
de las igualdades, sin alterar los resultados previos de acción, área,
capacidad y transporte.
% END THERMAL_R03_09
